\begin{table}[!htb]
\centering
\begin{threeparttable}
\caption{Pension system reform, year 2010.}
\label{table:Argentina:Universal}
\begin{tabular}{lccc}
\toprule
\textbf{Scenario} & \textbf{Tax rate} & \textbf{Savings rate} & \textbf{Avg. workweek (hours)}\\
\midrule 
Baseline & 12.5\% & 18.4\% & 42.54 \\[1ex]
Economic Equilibrium & 12.5\% & 18.78\% & 42.48 \\[1ex]
Full effect & 15.56\% & 17.62\% & 42.16 \\[1ex]
% OLD RESULTS:
% \textit{Pre}-reform & 12.5\% & 18.4\% & 42.54 \\[1ex]
% Economic Equilibrium & 12.5\% & 18.78\% & 42.48 \\[1ex]
% \textit{Post}-reform & 15.92\% & 17.51\% & 42.11 \\[1ex]
\bottomrule
\end{tabular}
\begin{tablenotes}
\footnotesize
\item\textit{Note:} The \textit{post}-reform scenario is simulated by permanently shifting $\epsilon = 1-\theta + \theta \eta_1 h_1 / h$. The "economic equilibrium" scenarios is simulated by solving for the new economic equilibrium path with taxes as in the Baseline scenario, but with the new permanent $\epsilon$.   
\end{tablenotes}
\end{threeparttable}
\end{table}
